\begin{exercises}
    \subsection{Bài tập}

    \begin{questions}[resume=SGK]
        \item Trong các hình phẳng sau (H.9.52), hình nào là hình phẳng có dạng đa giác đều?

        % TODO(HINH-NGUON): bo sung asset/crop dung Hinh 9.52 tu SGK trang 89; khong redraw khi chua kiem chung chinh xac.

        \answerline

        \item Trong các hình dưới đây (H.9.53), hình nào vẽ hai điểm $M$ và $N$ thoả mãn phép quay thuận chiều $60^\circ$ tâm $O$ biến điểm $M$ thành điểm $N$?

        % TODO(HINH-NGUON): bo sung asset/crop dung Hinh 9.53 tu SGK trang 89; khong redraw khi chua kiem chung chinh xac.

        \answerline

        \item Cho tam giác đều $ABC$ nội tiếp đường tròn $(O)$ bán kính 2 cm. Tính độ dài các cạnh của tam giác $ABC$.

        \answerline
        \answerline

        \item Cho hình thoi $ABCD$ có $\widehat A=60^\circ$. Gọi $M$, $N$, $P$, $Q$ lần lượt là trung điểm của $AB$, $BC$, $CD$, $DA$. Chứng minh rằng $MBNPDQ$ là lục giác đều.

        \answerline
        \answerline
        \answerline

        \item Cho tam giác đều $ABC$ nội tiếp đường tròn $(O)$ như Hình 9.54. Phép quay ngược chiều $60^\circ$ tâm $O$ biến các điểm $A$, $B$, $C$ lần lượt thành các điểm $D$, $E$, $F$. Chứng minh rằng $ADBECF$ là một lục giác đều.

        \begin{center}
            \begin{tikzpicture}[scale=0.8]
                \coordinate (O) at (0,0);
                \coordinate (A) at (90:2);
                \coordinate (D) at (150:2);
                \coordinate (B) at (210:2);
                \coordinate (E) at (270:2);
                \coordinate (C) at (330:2);
                \coordinate (F) at (30:2);
                \draw (O) circle (2);
                \draw (A)--(B)--(C)--cycle;
                \fill (O) circle (1.2pt);
                \node[above=1pt of A] {$A$};
                \node[above left=1pt of D] {$D$};
                \node[left=1pt of B] {$B$};
                \node[below=1pt of E] {$E$};
                \node[right=1pt of C] {$C$};
                \node[above right=1pt of F] {$F$};
                \node[below left=1pt of O] {$O$};
            \end{tikzpicture}

            \textit{Hình 9.54}
        \end{center}

        \answerline
        \answerline
        \answerline

        \item Liệt kê năm phép quay giữ nguyên một ngũ giác đều nội tiếp một đường tròn tâm $O$.

        \answerline

        \item Cho vòng quay mặt trời gồm tám cabin như Hình 9.55. Hỏi để cabin $A$ di chuyển đến vị trí cao nhất thì vòng quay phải quay thuận chiều kim đồng hồ quanh tâm bao nhiêu độ?

        % TODO(HINH-NGUON): Hinh 9.55 la hinh minh hoa vong quay mat troi SGK trang 89; bo sung asset/crop dung nguon, khong redraw xap xi.

        \answerline
    \end{questions}
\end{exercises}
